\prefacesection{Lay Abstract}

A lay abstract of not more 150 words must be included explaining the key goals and contributions of the thesis in lay terms that is accessible to the general public. 
\\~\\
The Navier-Stokes equations are a set of equations used to model the behaviour of viscous incompressible fluids. The question of whether or not the Navier-Stokes equations, or the inviscid Euler equations, are a valid description of fluid mechanics remains unanswered. A singularity is the occurrence of an infinite derivative of the velocity at a single point in the fluid flow. If the equations form a singularity from smooth initial conditions, then they are not valid descriptions of fluid mechanics. The Euler-Voigt equations are an inviscid regularization of the Euler equations which are known to not form singularities from smooth initial conditions. However, by decreasing the regularization parameter we approximate Euler flows. We examine the Euler-Voigt equations with a PDE-constrained optimization problem to identify initial conditions that produce 'extreme' behaviour in numerical simulations.  \#\# We present numerical evidence that CONCLUSION \#\#